\section{Die \LaTeX-Anbindung}\label{sec:latex}

Mathematische Texte entstehen in \LaTeX, und Sprachmodelle schreiben bevorzugt \LaTeX. Die naheliegende Stelle, an der eine Akte geführt
werden kann, ist deshalb das Dokument selbst: Jede Aussage wird dort deklariert, wo sie im Text steht, und der \TeX-Compiler prüft die
Regeln, während er das Dokument setzt. QAM stellt dafür drei Pakete bereit, die nur auf dem \LaTeX-Kern aufbauen (expl3, \LaTeX{} ab
2022-06) und von keinem Anwendungsprojekt abhängen.

\begin{center}
\begin{tabular}{@{}lp{0.68\textwidth}@{}}
\toprule
Paket & Aufgabe\\
\midrule
\texttt{qam.sty} & Akte führen, Regeln (R1)--(R4) beim Kompilieren prüfen, Notation und Übersicht setzen, Akte exportieren\\
\texttt{qam-physics.sty} & Formeln in gewöhnlicher Schreibweise unter einem deklarierten Interpretationsprofil registrieren\\
\texttt{qam-readout.sty} & erzeugte Belegfragmente der Ausführungsmaschine in das Dokument einbinden\\
\bottomrule
\end{tabular}
\end{center}

\subsection{Laden und Betriebsarten}

\begin{lstlisting}
\usepackage{qam}                   % Regelverstoß: Fehler, Compiler hält an
\usepackage[mode=warning]{qam}     % Regelverstoß: nur Warnung
\usepackage[export]{qam}           % zusätzlich \jobname.qam schreiben
\end{lstlisting}
Im Normalfall bricht ein Regelverstoß das Kompilieren ab. Das ist beabsichtigt: Ein Dokument mit einer Beförderung soll nicht als PDF
entstehen. Die Warnbetriebsart dient dem Einführen in einen bestehenden Text; die dabei exportierte Akte weist der Lean-Prüfer ab, so
dass der Verstoß nicht unbemerkt bleibt.

\subsection{Deklarationen}

\begin{center}
\small
\begin{tabular}{@{}p{0.46\textwidth}p{0.48\textwidth}@{}}
\toprule
Befehl & Wirkung\\
\midrule
\verb|\qamframe{F}| & Frame deklarieren\\
\verb|\qammap[sym]{g}{A}{B}| & Abbildung \(g\colon A\to B\) deklarieren\\
\verb|\qamcompose[sym]{k}{g}{h}| & Komposition zweier Abbildungen deklarieren\\
\verb|\qaminput[sym]{key}{F}{status}| & eingesetzten Wert registrieren\\
\verb|\qamclaim[sym]{key}{F}{status}| & Behauptung ohne Prämissen registrieren\\
\verb|\qamthesis[sym]{key}{F}{status}{Wortlaut}| & Ausgangsthese; der Wortlaut geht in den Hash ein\\
\verb|\qamderive[sym]{key}{status}{p1,p2,...}| & Herleitung aus Prämissen; Frame und Eingaben werden berechnet\\
\verb|\qamconvert[sym]{key}{status}{g}{p}| & Überführung von \(p\) über \(g\)\\
\verb|\qamcompare{a}{b}| & Vergleich, nur im selben Frame\\
\verb|\qamindependent{key}| & Unabhängigkeitserklärung\\
\bottomrule
\end{tabular}
\end{center}
Das optionale Argument ist das mathematische Symbol, unter dem der Eintrag im Text erscheint. Frame und Eingabemenge einer Herleitung
werden nicht angegeben, sondern aus den Prämissen berechnet; der Autor kann sie also nicht falsch angeben. Den Status gibt er an, und
das Paket prüft ihn gegen (R2). Leerzeichen in Prämissenlisten werden ignoriert.

Bei einem Verstoß meldet der Compiler die betroffene Regel mit den beteiligten Schlüsseln, zum Beispiel:
\begin{lstlisting}
! Package qam Error: Promotion: `beta' is declared `proved', but its weakest
(qam)                premise `gamma' is only `conditional'.
\end{lstlisting}
Entsprechende Meldungen gibt es für Frame-Vermischung, für eine Unabhängigkeitserklärung über einem Eintrag mit Eingaben sowie für
undeklarierte und doppelte Schlüssel, Frames und Abbildungen.

\subsection{Erzeugte Notation}

Die Deklarationen enthalten alles, was für eine kompakte Schreibweise nötig ist, und das Paket setzt sie selbst:
\begin{center}
\small
\begin{tabular}{@{}p{0.30\textwidth}p{0.64\textwidth}@{}}
\toprule
Befehl & Ausgabe\\
\midrule
\verb|\qam{key}| & Eintrag in QAM-Notation mit Status\\
\verb|\qam*{key}| & dasselbe ohne Status\\
\verb|\qamshow{key}| & Kurzform: Status, Frame, Eingaben\\
\verb|\qamhash{key}|, \verb|\qamhash*{key}| & kurzer bzw.\ voller Hash\\
\verb|\qamrecordhash| & Hash der ganzen Akte\\
\verb|\qamledger| & Tabelle aller Einträge in Deklarationsreihenfolge mit Hashes\\
\bottomrule
\end{tabular}
\end{center}

\paragraph{Ein Beispiel, das in diesem Dokument ausgewertet wird.} Der folgende Block steht so im Quelltext dieses Artikels; das Paket
\texttt{qam.sty} ist geladen und prüft ihn beim Setzen.
\begin{lstlisting}
\qamframe{M} \qamframe{T}
\qammap[G]{G}{M}{T}
\qaminput[\gamma]{gamma}{M}{computed}
\qamclaim[\kappa]{counting}{M}{proved}
\qamderive[\beta]{beta}{computed}{gamma,counting}
\qamconvert[\beta']{betaT}{conditional}{G}{beta}
\end{lstlisting}
\qamframe{M}\qamframe{T}
\qammap[G]{G}{M}{T}
\qaminput[\gamma]{gamma}{M}{computed}
\qamclaim[\kappa]{counting}{M}{proved}
\qamderive[\beta]{beta}{computed}{gamma,counting}
\qamconvert[\beta']{betaT}{conditional}{G}{beta}
Die Befehle \verb|\qam{gamma}|, \verb|\qam{beta}| und \verb|\qam{betaT}| setzen daraus
\begin{center}
\qam{gamma}\qquad \qam{beta}\qquad \qam{betaT}
\end{center}
und \verb|\qamshow{betaT}| ergibt \qamshow{betaT}. Man liest an der Notation ab, dass \(\beta'\) auf der Eingabe \(\gamma\) ruht, über \(G\)
aus dem Frame \(M\) gebracht wurde und höchstens bedingt gilt. Hätte der Quelltext \verb|\qamderive{beta}{proved}{gamma,counting}|
gelautet, wäre dieser Artikel nicht kompiliert worden. Die Übersicht \verb|\qamledger| für diese Einträge:
\begin{center}
\small
\qamledger
\end{center}

\subsection{Export}

Mit der Option \texttt{export} schreibt das Paket jeden Eintrag als Zeile in \texttt{\textbackslash jobname.qam}. Das Format ist
zeilenorientiert und für Menschen lesbar:
\begin{lstlisting}
% QAM record exported by qam.sty
% source example.tex md5 31EAB103...
map G : M -> T
gamma := <6,gamma>; frame M; status computed; inputs {gamma}; hash C4A2B3DB...; sources {}
counting := claim; frame M; status proved; inputs {}; hash 8097F8AB...; sources {}
beta := derive(gamma,counting); frame M; status computed; inputs {gamma};
        hash B0541AFD...; sources {C4A2B3DB...,8097F8AB...}
betaT := convert(G,beta); frame T; status conditional; inputs {gamma}; ...
% record C495463D...
\end{lstlisting}
(Hashes gekürzt, die Zeile von \texttt{beta} umbrochen.) Eine Eingabe erscheint als Zeugenobjekt \(\langle6,\cdot\rangle\) der
Objektebene. Die exportierte Datei ist die Schnittstelle zu allen weiteren Prüfern: Sie wird von Lean geprüft, von einem Skript
nachgerechnet und von der Datenbank indiziert (Abschnitt~\ref{sec:pruefer}). Damit hängt die Prüfung nicht davon ab, dass man dem
\TeX-Paket vertraut.

\subsection{Formeln unter einem Interpretationsprofil}\label{sec:profile}

Eine Formel wie die Einsteingleichung oder die Schrödingergleichung bedeutet nur in einem bestimmten Rahmen etwas. Das Paket
\texttt{qam-physics.sty} macht diesen Rahmen zu einem registrierten Objekt:
\begin{lstlisting}
\qamframe{ART}\qamframe{QM}
\qamprofile{gr}{ART}{L_{\rm ART}}{conditional}
  {Classical GR with Lorentz metric and Levi-Civita connection, SI units}
\begin{qamcontext}{gr}
\[ \qamformula{einstein}{conditional}
     {G_{\mu\nu}+\Lambda g_{\mu\nu}=\frac{8\pi G_{\rm N}}{c^4}T_{\mu\nu}} \]
\end{qamcontext}
\end{lstlisting}
Ein Profil (\verb|\qamprofile{key}{frame}{sprache}{status}{interpretation}|) ist unveränderlich und als gehashte Prämisse Teil jeder unter ihm
registrierten Formel: Wer die Interpretation ändert, ändert den Hash der Formel. Innerhalb von \verb|qamcontext| erscheinen Formeln in
gewöhnlicher Schreibweise. \verb|\qamphys{einstein}| verwendet die gespeicherte Formel wieder und zeigt außerhalb ihres Profils
automatisch Frame und Sprache an; \verb|\qamphys*| zeigt sie immer. Kontexte lassen sich verschachteln, das äußere Profil wird danach
wiederhergestellt. Ein Übergang in ein anderes Profil ist nur ausdrücklich möglich,
\begin{lstlisting}
\qamphysconvert{key}{status}{abbildung}{quelle}{zielprofil}{zielformel}
\end{lstlisting}
mit deklarierter Abbildung und vom Autor angegebener Zielformel. Auch zwei Profile im selben Frame werden nicht stillschweigend
gemischt. Das Paket prüft Kontext, Profilidentität, Quell- und Ziel-Frame, Doppeldeklarationen und vererbten Status. Die Formel selbst
ist ein unexpandierter \TeX-Text: Tensoren, Einheiten, Definitionsbereiche und physikalische Gesetze werden nicht geprüft, und eine
Übersetzung hebt den Status nie an.

\subsection{Belege der Ausführungsmaschine im Dokument}

Ergebnisse der Ausführungsmaschine (Abschnitt~\ref{sec:maschine}) werden nicht abgetippt, sondern als erzeugte Fragmente eingebunden.
Vor dem Kompilieren schreibt die Maschine ein Fragment,
\begin{lstlisting}
python -m qam.tools.readout_db mma KONTEXT.json AUFTRAG.json --tex-output build/ergebnis.tex
\end{lstlisting}
und das Dokument bindet es mit \verb|\usepackage{qam-readout}| und \verb|\qamreadout{build/ergebnis.tex}| ein. Das Fragment zeigt den
Ausgang, das Ergebnis oder die Gründe für eine Sperre und den SHA-512-Bezeichner des Belegs. Fehlt das Fragment, erscheint eine
sichtbare Platzhaltermeldung und eine Warnung. \LaTeX{} ruft dabei weder Python noch die Datenbank auf. Umgekehrt stellt
\texttt{receipt HASH --tex-output ...} aus einem im Dokument gedruckten Hash den gespeicherten Auftrag und Beleg wieder her.

\subsection{Übersetzen und Testen}

Die Pakete liegen in einem Verzeichnis, das in den Suchpfad von \TeX{} aufgenommen wird, etwa
\begin{lstlisting}
TEXINPUTS=pfad/zu/qam/tex: pdflatex dokument.tex
\end{lstlisting}
Eine plattformunabhängige Testreihe übersetzt gültige und absichtlich fehlerhafte Beispieldokumente, prüft, dass die gültigen
übersetzen und die fehlerhaften anhalten, und kontrolliert, falls \texttt{pdftotext} verfügbar ist, die gesetzte Ausgabe:
\begin{lstlisting}
python3 qam/tex/tests/run_tests.py --output-dir build/qam-tests
\end{lstlisting}
Das sind Softwaretests des Pakets, keine Prüfungen der registrierten Inhalte.
